% ============================================================================
% decomposer: A Model-Agnostic Framework for Factorial Decomposition
%
% This manuscript describes the decomposer R package and proves that
% Das Gupta's symmetric standardization is a special case of the
% Shapley value decomposition.
% ============================================================================
\documentclass[11pt]{article}

% ---- Packages ----
\usepackage[utf8]{inputenc}
\usepackage[T1]{fontenc}
\usepackage{amsmath, amssymb, amsthm}
\usepackage{booktabs}
\usepackage{geometry}
\usepackage{graphicx}
\usepackage{hyperref}
\usepackage{natbib}
\usepackage{algorithm}
\usepackage{algpseudocode}
\usepackage{listings}
\usepackage{xcolor}
\usepackage{enumitem}

\geometry{margin=1in}

% ---- Theorem environments ----
\newtheorem{theorem}{Theorem}
\newtheorem{proposition}[theorem]{Proposition}
\newtheorem{corollary}[theorem]{Corollary}
\newtheorem{definition}[theorem]{Definition}
\newtheorem{remark}[theorem]{Remark}

% ---- Code listing style ----
\lstset{
  language=R,
  basicstyle=\ttfamily\small,
  keywordstyle=\color{blue},
  commentstyle=\color{gray},
  stringstyle=\color{red!70!black},
  showstringspaces=false,
  breaklines=true,
  frame=single,
  numbers=left,
  numberstyle=\tiny\color{gray},
  xleftmargin=2em,
  framexleftmargin=1.5em,
  captionpos=b
}

% ---- Macros ----
\newcommand{\pkg}[1]{\texttt{#1}}
\newcommand{\code}[1]{\texttt{#1}}
\newcommand{\R}{\textsf{R}}

\title{%
  \pkg{decomposer}: A Model-Agnostic Framework for Factorial Decomposition
  of Discrete-Time Simulation Outcomes \\[0.5em]
  \large With a Proof of Equivalence Between Das~Gupta and
  Shapley Value Decomposition}

\author{
  Austin Carter \\
  Institute for Health Metrics and Evaluation \\
  University of Washington \\
  \texttt{aucarter@uw.edu}
}

\date{\today}

\begin{document}

\maketitle

% ====================================================================
\begin{abstract}
Factorial decomposition methods are widely used in demography,
epidemiology, and health economics to attribute changes in aggregate
outcomes to individual driving factors.
Two prominent approaches---Das~Gupta's symmetric standardization from
demography and Shapley value decomposition from cooperative game
theory---are often treated as distinct methods in the applied literature.
We prove that the two are algebraically identical for factorial
decomposition problems: the Das~Gupta weights are a special case of the
Shapley value weights when the characteristic function is defined over
binary factor coalitions.
Building on this unification, we present \pkg{decomposer}, an \R{}
package that provides a generic, model-agnostic decomposition framework
for arbitrary calibrated discrete-time simulation models.
The package defines an S3 interface consisting of four required methods
and four optional methods; any simulation model that implements this
interface can be decomposed without modifying the core algorithm.
We describe the mathematical framework, software design, and
computational optimizations (including warm-start support), and
illustrate the package with an SIR epidemiological model.
\end{abstract}

\bigskip
\noindent\textbf{Keywords:} decomposition, Shapley value, Das~Gupta,
counterfactual analysis, R package, simulation


% ====================================================================
\section{Introduction}
\label{sec:introduction}

Understanding the relative contributions of multiple factors to an
observed change is a fundamental analytical challenge.
In demography, Das~Gupta developed a symmetric standardization approach
for decomposing differences in rates into contributions from
compositional and rate-specific components
\citep{dasgupta1993standardization}.
Independently, game-theoretic decomposition based on Shapley values
\citep{shapley1953value} has been adopted in economics, political
science, and, more recently, in global health and epidemiological
modeling \citep{carter2024pediatric, shorrocks2013decomposition}.
Both approaches share desirable properties: they produce a unique,
symmetric decomposition in which the factor contributions sum exactly to
the total observed change.

Despite their shared properties, the two methods are typically presented,
implemented, and discussed as though they were different techniques.
Applied researchers may choose one framework over the other based on
disciplinary convention without recognizing their mathematical
equivalence.
In Section~\ref{sec:equivalence} we prove that, for the class of
factorial decomposition problems (where each factor can be switched
between an observed and a counterfactual value), the Das~Gupta weights
and Shapley value weights are identical.

On the software side, decomposition analyses are frequently embedded
within model-specific code, making them difficult to reuse, validate,
and extend.
Each new simulation model requires a bespoke implementation of the
decomposition procedure, despite the underlying mathematical algorithm
being the same.
To address this, we present \pkg{decomposer}, an \R{} package that
separates the decomposition algebra from the simulation model through a
generic S3 interface.
Any calibrated discrete-time simulation model---whether an SIR
compartmental model, an HIV transmission model operating via matrix
projection, or a demographic microsimulation---can use the
\pkg{decomposer} framework by implementing four S3 methods that define
how to run the model, summarize its outputs, construct counterfactual
inputs, and retrieve baseline values.

The remainder of this paper is organized as follows.
Section~\ref{sec:framework} presents the mathematical framework for
factorial decomposition.
Section~\ref{sec:equivalence} proves the equivalence of Das~Gupta and
Shapley decomposition.
Section~\ref{sec:software} describes the software design and interface.
Section~\ref{sec:example} provides a worked example.
Section~\ref{sec:discussion} concludes with a discussion.


% ====================================================================
\section{Mathematical Framework}
\label{sec:framework}

\subsection{Setup and notation}

Consider a discrete-time simulation model whose outcome of interest at
time $t$ depends on $n$ factors
$\mathbf{x} = (x_1, x_2, \ldots, x_n)$:
\[
  Y_t = f_t(x_1, x_2, \ldots, x_n).
\]
Each factor $x_i$ can take an \emph{observed} value $x_i^1$ (the actual
value used in the calibrated model) or a \emph{counterfactual} value
$x_i^0$ (the value in the absence of the factor, or its value at a
reference year).
In practice, the factors are often time-varying inputs such as treatment
coverage, transmission rates, or demographic parameters, and each
factor's observed and counterfactual trajectories span the full
simulation period.

The total change to be decomposed is:
\begin{equation}
\label{eq:total-change}
  \Delta Y_t = f_t(x_1^1, \ldots, x_n^1)
             - f_t(x_1^0, \ldots, x_n^0),
\end{equation}
and the goal is to write $\Delta Y_t = \sum_{i=1}^{n} \delta_i$ where
$\delta_i$ represents the contribution of factor $i$.

\subsection{Factorial design}

For each of the $2^n$ binary vectors
$\mathbf{b} = (b_1, \ldots, b_n) \in \{0, 1\}^n$, define the scenario
in which factor $i$ takes value $x_i^{b_i}$:
\[
  Y_t(\mathbf{b}) = f_t\!\bigl(x_1^{b_1}, \ldots, x_n^{b_n}\bigr).
\]
The \emph{marginal effect} of switching factor $i$ from counterfactual
to observed, holding the other factors at the configuration specified by
$S \subseteq N \setminus \{i\}$ (i.e., factors in $S$ are observed;
factors in $N \setminus S \setminus \{i\}$ are counterfactual), is:
\begin{equation}
\label{eq:marginal}
  m_i(S)
  = f_t(x_i^1, \, x_S^1, \, x_{N \setminus S \setminus \{i\}}^0)
  - f_t(x_i^0, \, x_S^1, \, x_{N \setminus S \setminus \{i\}}^0),
\end{equation}
where $N = \{1, \ldots, n\}$ and the notation $x_S^1$ means
$\{x_j^1 : j \in S\}$.

Both the Das~Gupta and Shapley methods define $\delta_i$ as a weighted
average of these marginal effects over all subsets
$S \subseteq N \setminus \{i\}$.


% ====================================================================
\subsection{Das~Gupta's symmetric standardization}
\label{sec:dasgupta}

Das~Gupta's method \citep{dasgupta1993standardization} computes the
contribution of factor $i$ as a symmetrically weighted average of
marginal effects, where the weight depends only on the number of other
factors in the ``on'' state:
\begin{equation}
\label{eq:dasgupta}
  \delta_i^{\text{DG}}
  = \sum_{\substack{S \subseteq N \setminus \{i\}}}
    w_{|S|}^{\text{DG}} \cdot m_i(S),
  \qquad
  w_k^{\text{DG}} = \frac{1}{n \binom{n-1}{k}},
  \quad k = 0, 1, \ldots, n-1.
\end{equation}
The weights are normalized so that the decomposition is exact, i.e.,
$\sum_{i=1}^n \delta_i^{\text{DG}} = \Delta Y_t$.
Das~Gupta derived these weights from the principle of symmetric
treatment of factors: no factor should receive systematically higher
weight by virtue of the order in which it is considered.


% ====================================================================
\subsection{Shapley value decomposition}
\label{sec:shapley}

The Shapley value \citep{shapley1953value} originates in cooperative
game theory, where it provides the unique allocation of a coalition's
payoff that satisfies the axioms of efficiency, symmetry, additivity,
and the null-player property.

To apply the Shapley value to factorial decomposition, define a
\emph{characteristic function} $v : 2^N \to \mathbb{R}$ by:
\begin{equation}
\label{eq:characteristic}
  v(S) = f_t(x_S^1, \, x_{N \setminus S}^0),
  \qquad S \subseteq N.
\end{equation}
This function maps each coalition of factors (those taking their
observed values) to the simulation outcome.
The Shapley value for factor $i$ is:
\begin{equation}
\label{eq:shapley}
  \phi_i
  = \sum_{S \subseteq N \setminus \{i\}}
    \frac{|S|! \, (n - |S| - 1)!}{n!}
    \bigl[ v(S \cup \{i\}) - v(S) \bigr].
\end{equation}
Observe that $v(S \cup \{i\}) - v(S) = m_i(S)$ exactly as in
Equation~\eqref{eq:marginal}, so the Shapley value is also a weighted
average of marginal effects:
\begin{equation}
\label{eq:shapley-weight}
  \phi_i
  = \sum_{S \subseteq N \setminus \{i\}}
    w_{|S|}^{\text{Sh}} \cdot m_i(S),
  \qquad
  w_k^{\text{Sh}} = \frac{k! \, (n - k - 1)!}{n!}.
\end{equation}


% ====================================================================
\section{Equivalence of Das~Gupta and Shapley Decomposition}
\label{sec:equivalence}

\begin{theorem}
\label{thm:equivalence}
For the class of factorial decomposition problems defined in
Section~\ref{sec:framework}, the Das~Gupta weight $w_k^{\mathrm{DG}}$
and the Shapley weight $w_k^{\mathrm{Sh}}$ are identical for all
$k = 0, 1, \ldots, n-1$.  That is:
\[
  \frac{k! \, (n - k - 1)!}{n!}
  = \frac{1}{n \binom{n-1}{k}}.
\]
Consequently, $\delta_i^{\mathrm{DG}} = \phi_i$ for every factor $i$.
\end{theorem}

\begin{proof}
We expand the right-hand side.  Recall that
$\binom{n-1}{k} = \frac{(n-1)!}{k!\,(n-1-k)!}$, so:
\begin{align*}
  \frac{1}{n \binom{n-1}{k}}
  &= \frac{1}{n \cdot \dfrac{(n-1)!}{k!\,(n-1-k)!}}
  \\[6pt]
  &= \frac{k! \, (n - 1 - k)!}{n \cdot (n-1)!}
  \\[6pt]
  &= \frac{k! \, (n - 1 - k)!}{n!},
\end{align*}
where the last step uses $n \cdot (n-1)! = n!$.
Since $n - k - 1 = n - 1 - k$, both sides are:
\[
  w_k^{\text{DG}} = w_k^{\text{Sh}} = \frac{k!\,(n-k-1)!}{n!}.
\]
Thus the weight assigned to each marginal effect $m_i(S)$ is the same
under both methods, and the decomposition results are identical.
\end{proof}

\begin{remark}
This result shows that Das~Gupta's approach and the Shapley value
decomposition, despite arising from different intellectual
traditions---demographic standardization and cooperative game
theory, respectively---yield the same decomposition for any factorial
design with binary factors.
Practitioners can therefore use whichever framing is most natural for
their discipline while obtaining identical quantitative results.
\end{remark}

\begin{remark}
The weights can also be expressed as:
\[
  w_k = \frac{1}{n \binom{n-1}{k}}
\]
which has an intuitive interpretation: there are $\binom{n-1}{k}$
subsets of size $k$ among the $n-1$ other factors; averaging uniformly
over all $n-1$ possible subset sizes requires dividing by $n$ (one for
each possible subset size $k = 0, \ldots, n-1$) and then by the number
of subsets of that size.
\end{remark}

\begin{corollary}
The \pkg{decomposer} package implements a single weight formula
$w_k = \bigl[n \binom{n-1}{k}\bigr]^{-1}$ that serves both the
\code{das\_gupta} and \code{shapley} method options.
The two options are retained for user-facing clarity, not because they
produce different results.
\end{corollary}


% ====================================================================
\section{Counterfactual Types}
\label{sec:counterfactual}

The interpretation of the decomposition depends critically on the
definition of the counterfactual values $x_i^0$.
The \pkg{decomposer} package supports two counterfactual types, each
corresponding to a different analytical question.

\subsection{Change-since counterfactual}

Under the \emph{change-since} counterfactual (\code{cf\_type =
"change\_since"}), the counterfactual value $x_i^0$ is the factor's
value at a fixed baseline year $t_0$.
For all years $t > t_0$, the factor is held at its baseline-year level
rather than evolving to its observed trajectory.
This measures the impact of \emph{changes in the factor since the
baseline year}:
\[
  x_i^0(t) = x_i(t_0), \qquad \forall\, t \geq t_0.
\]
This is appropriate when the research question is:
\emph{How much of the change in the outcome since year $t_0$ is
attributable to changes in each factor?}

\subsection{Total-impact counterfactual}

Under the \emph{total-impact} counterfactual (\code{cf\_type =
"total\_impact"}), the counterfactual value $x_i^0$ is zero (or
whatever the model defines as the absence of the factor):
\[
  x_i^0(t) = 0, \qquad \forall\, t.
\]
This measures the \emph{total impact} of the factor, answering:
\emph{How many additional (or fewer) outcomes would occur if this factor
had never been present?}


% ====================================================================
\section{Software Design}
\label{sec:software}

\subsection{Overview}

The \pkg{decomposer} package separates the decomposition algorithm from
the simulation model through an S3 generic interface (Figure~\ref{fig:architecture}).
The core package implements the factorial run-table generation,
orchestration of counterfactual simulation runs, and the decomposition
algebra.
The user supplies a model object with methods that define the
model-specific behavior: running simulations, summarizing outputs, and
constructing counterfactuals.

\begin{figure}[ht]
\centering
\fbox{\parbox{0.85\textwidth}{\ttfamily
\begin{tabular}{l}
\textbf{decomposer (core package)} \\
\quad decompose() \textrm{--- orchestrator} \\
\quad make\_run\_table() \textrm{--- $2^n$ factorial design} \\
\quad run\_decomp() / decomp\_algebra() \textrm{--- weight calculation} \\[6pt]
\textbf{S3 interface (user implements)} \\
\quad run\_simulation.\emph{model}() \textrm{--- run the model} \\
\quad summarize\_results.\emph{model}() \textrm{--- extract outcomes} \\
\quad build\_counterfactual.\emph{model}() \textrm{--- construct CF params} \\
\quad get\_baseline\_value.\emph{model}() \textrm{--- baseline reference} \\[6pt]
\textbf{Optional methods} \\
\quad prepare\_decomp.\emph{model}() \textrm{--- one-time setup} \\
\quad supports\_warm\_start.\emph{model}() \textrm{--- warm-start flag} \\
\quad get\_warm\_start\_state.\emph{model}() \\
\quad run\_from\_warm\_start.\emph{model}()
\end{tabular}
}}
\caption{Architecture of the \pkg{decomposer} package.  The core package
  handles decomposition logic; users implement S3 methods for their
  specific model.}
\label{fig:architecture}
\end{figure}

\subsection{Required interface methods}

Any model that wishes to use the \pkg{decomposer} framework must
implement four S3 methods for its class:

\begin{description}[leftmargin=1.5em, labelindent=0em]
\item[\code{run\_simulation(model, params, ...)}]
  Execute the simulation model with the given (possibly counterfactual)
  parameters.  Returns raw simulation output in any model-specific
  format.

\item[\code{summarize\_results(model, sim\_output, ...)}]
  Transform raw simulation output into a \code{data.table} with a
  \code{year} column and one column per summary variable (e.g.,
  \code{deaths}, \code{infections}).

\item[\code{build\_counterfactual(model, params, factor\_name,
  decomp\_year, cf\_type, baseline\_values, ...)}]
  Given observed parameters, produce a modified parameter set where the
  specified factor is replaced by its counterfactual value from
  \code{decomp\_year} onward.

\item[\code{get\_baseline\_value(model, params, factor\_name, cf\_type,
  baseline\_year, ...)}]
  Return the baseline value for a factor.  For \code{change\_since}
  counterfactuals, this is the factor's value at the baseline year;
  for \code{total\_impact}, it is the zero/null value.
\end{description}

\subsection{Optional methods for performance}

For large-scale models, re-running the simulation from the beginning of
time for each of the $2^n$ counterfactual combinations at each
decomposition year can be prohibitively expensive.
The \pkg{decomposer} package provides an optional \emph{warm-start}
mechanism:

\begin{description}[leftmargin=1.5em, labelindent=0em]
\item[\code{prepare\_decomp(model, params, ...)}]
  One-time preparation before the decomposition loop.  The default
  implementation runs a full observed simulation and caches the result.

\item[\code{supports\_warm\_start(model, ...)}]
  Returns \code{TRUE} if the model supports warm starts.

\item[\code{get\_warm\_start\_state(model, sim\_output, year, ...)}]
  Extracts the model state at a given year from a full simulation,
  enabling restart without re-simulating prior years.

\item[\code{run\_from\_warm\_start(model, warm\_state, params,
  output\_years, ...)}]
  Runs the simulation forward from a saved state with (possibly
  counterfactual) parameters, producing output only for the
  requested years.
\end{description}

With warm starts, the cost of decomposing $T$ years with $n$ factors is
reduced from $O(T \cdot 2^n \cdot C_{\text{full}})$ to
$O(T \cdot 2^n \cdot C_{\text{partial}})$ where
$C_{\text{partial}} \ll C_{\text{full}}$ because only the years from
the decomposition year onward need to be simulated.

\subsection{Orchestration pipeline}

Algorithm~\ref{alg:decompose} summarizes the main \code{decompose()}
function, which orchestrates the full pipeline.

\begin{algorithm}[ht]
\caption{The \code{decompose()} orchestration pipeline.}
\label{alg:decompose}
\begin{algorithmic}[1]
\Require Model object \code{model}, parameters \code{params}, factors
  $\{x_1, \ldots, x_n\}$, summary variables, decomposition years
  $\{t_1, \ldots, t_T\}$
\Ensure Decomposition results (standard, step scenarios, extended)
\State \textbf{Prepare:}
  \code{prepare\_decomp(model, params)} $\to$
  cached full simulation and summary
\State \textbf{Baseline values:}
  $\forall\, i$: \code{get\_baseline\_value(model, params, $x_i$,
  cf\_type, $t_0$)}
\State \textbf{Run table:}
  Generate $2^n$ binary combinations via
  \code{make\_run\_table(factors)}
\For{each decomposition year $t \in \{t_1, \ldots, t_T\}$}
  \For{each row $\mathbf{b}$ in the run table}
    \If{$\mathbf{b} = (1,\ldots,1)$}
      \State Use cached full simulation summary
    \Else
      \State Build counterfactual params for factors where $b_i = 0$
      \If{warm start supported}
        \State \code{run\_from\_warm\_start(model, full\_sim,
          params\_cf, $t$)}
      \Else
        \State \code{run\_simulation(model, params\_cf)}
      \EndIf
      \State \code{summarize\_results(model, sim\_output)}
    \EndIf
  \EndFor
  \State Apply decomposition algebra
    (Equation~\ref{eq:dasgupta} / \ref{eq:shapley-weight})
\EndFor
\State \Return standard decomp, step scenarios, extended decomp
\end{algorithmic}
\end{algorithm}


\subsection{Factorial run table}

The \code{make\_run\_table()} function generates the $2^n$ factorial
design as a \code{data.table} with one row per combination and one
column per factor indicating whether it is at its observed value (1)
or counterfactual value (0):

\begin{lstlisting}[caption={Run table for $n=3$ factors.},
                   label=lst:runtable]
make_run_table(c("A", "B", "C"))
#    run_id A B C
# 1:      1 0 0 0
# 2:      2 1 0 0
# 3:      3 0 1 0
# 4:      4 1 1 0
# 5:      5 0 0 1
# 6:      6 1 0 1
# 7:      7 0 1 1
# 8:      8 1 1 1
\end{lstlisting}


% ====================================================================
\section{Example: SIR Model Decomposition}
\label{sec:example}

To illustrate the \pkg{decomposer} framework, we implement a simple
discrete-time SIR (Susceptible--Infected--Recovered) compartmental model
with two decomposable intervention factors: (1) a reduction in the
transmission rate $\beta$ and (2) an increase in the recovery rate
$\gamma$ from treatment.

\subsection{Model specification}

The SIR model evolves a population of size $N$ through yearly time
steps:
\begin{align}
  S_{t+1} &= S_t - \beta_t \, S_t \, I_t / N, \label{eq:sir-s} \\
  I_{t+1} &= I_t + \beta_t \, S_t \, I_t / N
             - \gamma_t \, I_t - \mu \, I_t, \label{eq:sir-i} \\
  R_{t+1} &= R_t + \gamma_t \, I_t, \label{eq:sir-r}
\end{align}
where $\beta_t = \beta_0 (1 - r_t)$ is the effective transmission rate
after applying a coverage-dependent reduction $r_t$,
$\gamma_t = \gamma_0 + \tau_t$ is the effective recovery rate with
additional treatment effect $\tau_t$, and $\mu = 0.01$ is the disease
mortality rate.

The two factors to decompose are:
\begin{itemize}
  \item \textbf{Transmission reduction} ($r_t$): gradually increases
    from 0 to 0.5 over the simulation period.
  \item \textbf{Treatment effect} ($\tau_t$): gradually increases from
    0 to 0.15 over the simulation period.
\end{itemize}

\subsection{Interface implementation}

Implementing the four required methods for the \code{sir\_model} class
is straightforward:

\begin{lstlisting}[caption={Implementing the \pkg{decomposer} interface for the SIR model.},
                   label=lst:sir]
# 1. Run simulation
run_simulation.sir_model <- function(model, params, ...) {
  # Discrete-time SIR loop (Equations 6-8)
  # Returns data.table(year, S, I, R, infections, deaths)
}

# 2. Summarize results
summarize_results.sir_model <- function(model, sim_output, ...) {
  sim_output[, .(year, infections, deaths)]
}

# 3. Get baseline value
get_baseline_value.sir_model <- function(model, params,
    factor_name, cf_type, baseline_year, ...) {
  if (cf_type == CF_TOTAL_IMPACT) return(list(type="zero"))
  list(type = "change_since",
       value = params[[factor_name]][as.character(baseline_year)])
}

# 4. Build counterfactual
build_counterfactual.sir_model <- function(model, params,
    factor_name, decomp_year, cf_type, baseline_values, ...) {
  idx <- which(params$years >= decomp_year)
  if (cf_type == CF_TOTAL_IMPACT) {
    params[[factor_name]][idx] <- 0
  } else {
    params[[factor_name]][idx] <- baseline_values[[factor_name]]$value
  }
  params
}
\end{lstlisting}

\subsection{Running the decomposition}

With the interface implemented, running the decomposition requires a
single call:

\begin{lstlisting}[caption={Running the decomposition.}, label=lst:run]
library(decomposer)

model  <- create_sir_model(years = 1990:2020)
params <- create_sir_params(years = 1990:2020)

results <- decompose(
  model        = model,
  params       = params,
  factors      = c("beta_reduction", "treatment_effect"),
  summary_vars = c("deaths", "infections"),
  decomp_years = 2000:2020,
  method       = "das_gupta",
  cf_type      = "change_since",
  baseline_year = 2000
)
\end{lstlisting}

The returned \code{results} object contains:
\begin{itemize}
  \item \code{results\$decomp} --- a \code{data.table} with one row per
    decomposition year per summary variable, with columns for each
    factor's contribution.
  \item \code{results\$step\_scenarios} --- time series of the observed
    outcome anchored at each decomposition year.
  \item \code{results\$extended\_decomp} --- factor contributions at
    all years beyond each decomposition year.
  \item \code{results\$run\_table} --- the $2^n$ factorial design used.
\end{itemize}


% ====================================================================
\section{Application: Pediatric HIV Decomposition}
\label{sec:application}

The \pkg{decomposer} package was developed in the context of an
analysis decomposing changes in pediatric HIV burden across 45 countries
in sub-Saharan Africa \citep{carter2024pediatric}.
Two simulation backends were used: a legacy SpectrumSoftware/simmod
matrix-projection model and the newer leapfrog model
\citep{eaton2021naomi}.
By implementing adapters for each backend
(\code{simmod\_model} and \code{leapfrog\_model} classes), the same
decomposition orchestration code was used for both, with the leapfrog
adapter additionally supporting warm-start optimization to reduce
computation time.
The analysis decomposed five factors---vertical transmission,
antiretroviral treatment, adult incidence trends, fertility, and
child mortality---across 24 annual decomposition years, requiring
$2^5 = 32$ simulation runs per location--year--draw combination.


% ====================================================================
\section{Discussion}
\label{sec:discussion}

\subsection{Unified perspective}

The proof that Das~Gupta and Shapley weights are identical
(Theorem~\ref{thm:equivalence}) offers a unified perspective on two
widely used decomposition methods.
For applied researchers, this means the choice between ``Das~Gupta
decomposition'' and ``Shapley value decomposition'' is one of framing
and interpretation rather than methodology.
Demographers accustomed to standardization-based reasoning and
economists versed in cooperative game theory are performing precisely
the same computation.

This equivalence also implies that the extensive axiomatic
characterizations of the Shapley value---efficiency, symmetry,
linearity, and the null-player property
\citep{shapley1953value}---automatically apply to the Das~Gupta method,
providing an additional theoretical foundation for a technique that was
originally motivated by practical considerations of symmetric treatment.

\subsection{Software contributions}

The \pkg{decomposer} package makes several practical contributions:

\begin{enumerate}
  \item \textbf{Separation of concerns.}
    The decomposition algebra is entirely decoupled from the simulation
    model.
    New models can be incorporated by implementing four methods without
    modifying the decomposition code.

  \item \textbf{Validation.}
    The \code{validate\_model\_interface()} function provides a
    smoke test that checks for the existence of required methods and
    runs a minimal simulation to verify the interface contract.

  \item \textbf{Performance.}
    The warm-start mechanism enables significant computational savings
    for models where the simulation state at an intermediate year can
    be extracted and used as an initial condition.

  \item \textbf{Flexibility.}
    Two counterfactual types (\code{change\_since} and
    \code{total\_impact}) support different analytical questions without
    requiring separate implementations.

  \item \textbf{Reproducibility.}
    Standardizing the decomposition procedure across models reduces the
    risk of implementation errors and facilitates code review.
\end{enumerate}

\subsection{Limitations and extensions}

The current implementation is limited to binary factorial designs where
each factor is either fully observed or fully counterfactual.
Extensions to partial factor levels (e.g., attributing effects to
changes in specific quantiles of a factor) or to continuous factor
spaces would require a generalized weighting scheme.
The $2^n$ computational complexity also limits the practical number of
factors, though this is inherent to full factorial decomposition and is
mitigated by the warm-start optimization.

Future work could extend the framework to support sequential
(time-ordered) decomposition as a complement to the current symmetric
approach, as well as uncertainty quantification through bootstrap or
posterior-draw-level decomposition.


% ====================================================================
\section{Conclusion}

We have proved that the Das~Gupta symmetric standardization and the
Shapley value decomposition produce identical results for factorial
decomposition problems, unifying two traditions in the decomposition
literature.
The \pkg{decomposer} R package implements this unified framework in a
model-agnostic fashion, enabling any calibrated discrete-time simulation
model to be decomposed by implementing a small S3 interface.
The package is available at
\url{https://github.com/aucarter/decomposer}.


% ====================================================================
% References
% ====================================================================
\bibliographystyle{plainnat}
\bibliography{references}

\end{document}
